\documentclass[10pt,a4paper]{amsart}
\usepackage[margin=19mm]{geometry}
\usepackage{amsmath,amssymb,mathtools}
\usepackage[scr=rsfs,cal=euler]{mathalfa}
\usepackage{xfrac}
\usepackage[unicode,hidelinks,bookmarks=false]{hyperref}
\newcommand{\ie}{i.e.\ }
\newcommand{\cf}{cf.\ }
\newcommand{\resp}{respectively\ }
\newcommand{\Subsection}{\subsection}
\providecommand{\nobreakdash}{\nobreak\hbox{-}}
\setlength{\emergencystretch}{2em}
\multlinegap0pt
\usepackage{bm}
\usepackage{bbm}
\usepackage{dsfont}
\usepackage{xspace}
\usepackage{cancel}
\usepackage{mathtools}
\usepackage{amsthm}
\makeatletter
\@ifundefined{theorem}{\newtheorem{theorem}{Theorem}}{}
\@ifundefined{lemma}{\newtheorem{lemma}[theorem]{Lemma}}{}
\makeatother
\title[Finite Sample Analyses for Continuous-time Linear Systems: S]{Finite Sample Analyses for Continuous-time Linear Systems: System Identification and Online Control}
\author{Hongyi Zhou, Jingwei Li, Jingzhao Zhang}
\date{Source: arXiv:2509.22741v1 (CC BY 4.0)}
\begin{document}
\maketitle

\noindent\emph{Source and licence.} Finite Sample Analyses for Continuous-time Linear Systems: System Identification and Online Control. Version arXiv:2509.22741v1 (CC BY 4.0), distributed under the Creative Commons Attribution 4.0 International licence, as recorded in the arXiv licence metadata for this exact version and shown on the abstract page https://arxiv.org/abs/2509.22741v1. Full rights notice: licenses/arxiv-2509.22741-CC-BY-4.0.txt.\par\smallskip

\noindent\textit{Editorial scope.} Counted unit: the Lemma whose source label is \texttt{lem:matrix inverse} in the article (source0.tex lines 539--549, source label \texttt{lem:matrix inverse}), delivered together with the article's own complete original proof of it (source0.tex lines 550--574, one \texttt{proof} environment of 25 lines). No mathematics is written, repaired or re-derived: statement and proof are the article's own text line for line. Required context retained uncounted: none. Every definition, display and label of those lines is retained unchanged. Omitted from this package and not redistributed: 1--538, 575--2094. Nothing inside a retained excerpt is omitted or altered: every retained body is delivered exactly as the article's lines are written, so no omission receipt of any kind accompanies this package. Citations: the retained excerpts carry no citation command at all, so no bibliography entry of the article is transcribed and no reference is redistributed. Reference adaptation: none was needed and none was made; every reference inside the delivered unit resolves inside it, so no unresolved reference or citation is produced. Printed slips kept literal: no printed word, symbol, hypothesis or number of the counted statement or of its proof was altered. Layout and class: the article's own class is \texttt{article} with packages including neurips\_2025, inputenc, fontenc, hyperref, url, booktabs, amsfonts, nicefrac, microtype, xcolor, graphicx, psfrag, amsmath, verbatim; the wrapper is the lane's pinned \texttt{amsart} editorial wrapper at 10pt on A4, which restates the theorem-like environments the retained text needs and adds the article's own macro definitions that the retained text needs (none), with extra wrapper packages (bm, bbm, dsfont, xspace, cancel, mathtools). The delivered numbering is the unit's own. Assets: no retained span contains a figure environment, a graphics-inclusion command, a PDF-page inclusion, a table or a TikZ picture; no image file is copied into this package, the package declares no asset dependency and no image file is redistributed with it. Licence: version arXiv:2509.22741v1 of this article is distributed under the Creative Commons Attribution 4.0 International licence, as recorded in the arXiv licence metadata for this exact version and shown on the abstract page https://arxiv.org/abs/2509.22741v1. A full rights notice is delivered at licenses/arxiv-2509.22741-CC-BY-4.0.txt. Characters: every retained excerpt is pure ASCII, so the delivered engine, XeLaTeX with its default Latin Modern text font, reports no missing character.
\begin{lemma}[Matrix inverse]\label{lem:matrix inverse}
For any $A\in\mathbb{R}^{d\times d}$ and $t$ such that $0<\Vert At\Vert \leq \frac{1}{10}$, we have the following estimation of $e^{At}$:
\begin{align*}
    \Vert e^{At}-I_{d}\Vert \leq e^{\Vert At\Vert }-1\,,
\end{align*}
and if we denote $A_{1}=e^{At}$, then $A$ also satisfies that 
\begin{align*}
A=\frac{1}{t}\sum_{k\geq 1}\frac{(-1)^{k+1}}{k}(A_{1}-I_{d})^{k}\,.
\end{align*}

\end{lemma}

\begin{proof}
We expand $e^{At}$ by
\begin{align*}
    e^{At}=\sum_{k\geq 0}\frac{1}{k!}(At)^{k}\,,
\end{align*}
which follows that
\begin{align*}
    \Vert e^{At}-I_{d}\Vert =\left\Vert\sum_{k\geq 1}\frac{1}{k!}(At)^{k}\right\Vert\leq \sum_{k\geq 1}\frac{1}{k!}\Vert At\Vert ^{k}= e^{\Vert At\Vert }-1\leq \frac{1}{9}\,.
\end{align*}

Since $\Vert A_{1}-I_{d}\Vert <1$, the progression $A_{2}=\sum_{k\geq 1}\frac{(-1)^{k+1}}{kt}(A_{1}-I_{d})^{k}$ converges, and thus 
$e^{A_{2}t}=e^{At}$. Furthermore, it can be computed that
\begin{align*}
\Vert A_{2}t\Vert \leq \sum_{k\geq 1}\left\Vert \frac{1}{k}(A_{1}-I_{d})\right\Vert\leq \sum_{k\geq 1}\frac{1}{k}(\frac{1}{9})^{k}\leq \frac{1}{8}\,.
\end{align*}

Now we show that $A_{2}=A$. We have already known that $\Vert At\Vert $ and $\Vert A_{2}t\Vert $ are small. We also note that the function $f:X\to e^{X}$ $\left( \Vert X\Vert \leq \frac{1}{8}\right)$ constitutes a one-to-one mapping. This assertion is supported by the observation that for any  $X_{1},X_{2}$ such that $\Vert X_{1}\Vert ,\Vert X_{1}+X_{2}\Vert \leq\frac{1}{8}$, we have $\Vert X_{2}\Vert \leq \frac{1}{4}$, implying that 
\begin{align}\label{error of exponential}
\left\Vert e^{X_{1}+X_{2}}-e^{X_{1}}-X_{2}\right\Vert&=\left\Vert\sum_{k\geq 2}\frac{1}{k!}(X_{1}+X_{2})^{k}-X_{1}^{k}\right\Vert\\
&\leq\sum_{k\geq 2}\frac{1}{k!}\frac{2^{k}-1}{4^{k-1}}\Vert X_{2}\Vert \\
&\leq \frac{1}{2}\Vert X_{2}\Vert \,.
\end{align}
Then $\left\Vert e^{X_{1}+X_{2}}-e^{X_{1}}\right\Vert \geq \frac{1}{2}\Vert X_{2}\Vert $, which means $f$ is one-to-one, and thereby leading that $A_{2}=A$.

\end{proof}

\end{document}
